\section{Introduction}\label{sec:introduction}

Multi-agent applications decompose a user task into specialized analyses
and dependent synthesis. A research report can combine independent study
of alternatives, application needs, and external constraints; a repository
review can combine behavior, implementation, and test analysis. Each Agent
can make several consecutive model calls, retaining its previous inputs and
outputs. These growing histories become worker-local KV state.

A serving-worker failure can remove the states of several Agents across
multiple jobs. Without a surviving compatible checkpoint, a replacement
worker must prefill the durable history before generation resumes
\cite{strati2024dejavu,cao2026lumen}. Full protection avoids that computation
but consumes memory and transfer resources during healthy serving. The
problem is therefore to spend a finite protection budget where it removes
the most workflow delay.

Three distinctions guide \system. First, reconstruction delay can be
absorbed by residual work on other dependency paths, so reconstruction cost
alone is insufficient to prioritize Agents. Second, complete checkpoint
coverage and fast availability serve different purposes. A complete
Peer-DRAM checkpoint avoids full reconstruction, while an additional
Peer-HBM subset only helps when it removes exposed transfer waiting.
Third, replica HBM competes with foreground execution and must track safe
capacity. The research-report and repository-review applications described in
Section~\ref{sec:background:serving} ground these questions; Section~\ref{sec:motivation} explains the
mechanisms.

\system's Tracker estimates remaining Agent-task duration from causally
available information and propagates guarded Slack through the residual
DAG. KV Master selects complete DRAM checkpoints using exposed
reconstruction delay, chooses causal Top-$K$ blocks in the first $L$
layers, and ranks HBM copies by their remaining exposed-delay benefit.
The Adapter adjusts the total HBM quota while preserving full DRAM
coverage. Router and Connector retain compatibility, failure-domain,
and committed-state requirements.

This manuscript presents the recovery design and a new evaluation
protocol. The protocol uses public task inputs and project-authored
applications, rather than relabeling existing traces as those applications.
It pairs failure-free and faulted executions at equal physical budgets,
measures directly interrupted and other co-served jobs, and isolates the
incremental benefit of HBM beyond complete DRAM coverage. The contributions
are:
\begin{itemize}
\item A workflow-aware protection policy connecting Agent delay tolerance
to complete checkpoint admission.
\item A fast-region policy combining causal native KV demand, layer-ordered
recovery, and safe HBM quota adaptation.
\item A reproducible experiment design for two multi-turn applications,
with matched baselines, common-failure analysis, component ablation,
quality checks, and explicit zero- or negative-benefit cases.
\end{itemize}
